\section{Součin rozpustnosti}

\subsection{Rozpustnost}
Udává kolik gramů látky se za dané teploty rozpustí ve 100 g rozpouštědla, zpravidla vody. Rozpustnost se zvyšující se teplotou zpravidla roste.

\begin{tabular}{|l|l|l|l|}
	\hline
	& \ce{NH4Cl} & NaCl & KCl \\\hline
	0 $^\circ$C & 29,41 & 35,63 & 28,14 \\\hline
	50 $^\circ$C & 50,42 & 36,70 & 43,05 \\\hline
	60 $^\circ$C & 55,30 & 37,08 & 45,88 \\\hline
\end{tabular}

\begin{itemize}
	\item \textbf{Nenasycený roztok} -- koncentrace látky je nižší, než její rozpustnost za dané teploty.
	\item \textbf{Nasycený roztok} -- koncentrace látky je shodná s její rozpustností za dané teploty.
	\item \textbf{Přesycený roztok} -- koncentrace látky je vyšší, než její rozpustnost za dané teploty.
\end{itemize}

\zadani{Kolik g KCl se vysráží, pokud 100 g nasyceného roztoku ochladíme z 60 na 0$^\circ$C}

Nejprve si spočítáme hmotnostní zlomky nasycených roztoků při jednotlivých teplotách:

$w_0 = \frac{28,14}{128,14} = 0,2196\\
w_{60} = \frac{45,88}{148,88} = 0,3145\\
m_{H2O,60} = 100 - 31,45 = 68,44$ g \ce{H2O}

Hmotnost nasyceného roztoku při 0$^\circ$C je $\frac{68,44}{0,7804}$, což je 87,70 g. Rozdíl oproti původním 100 g odpovídá hmotnosti vyloučeného KCl.

\odpoved{Po ochlazení se vyloučí 12,30 g KCl}

\subsection{Součin rozpustnosti}
Popisuje koncentrace iontů v nasycených roztocích.

\ce{AgNO3 + KCl -> AgCl v + KNO3}\\
\ce{AgCl <=> Ag+ + Cl-}

Rovnováhu v nasyceném roztoku chloridu stříbrného můžeme popsat pomocí rovnovážné konstanty:

$K = \frac{[\ce{Ag+}][\ce{Cl-}]}{[\textrm{AgCl}]}$

Z rovnovážné konstanty pak odvodíme \textit{součin rozpustnosti}:

$K_S = [\ce{Ag+}][\ce{Cl-}]\\
pK_S = -\log K_S$

\clearpage

\textbf{Hodnoty součinů rozpustnosti} najdeme v chemických tabulkách.

\begin{center}
\begin{tabular}{|l|r@{,}l|r@{,}l||l|r@{,}l|r@{,}l|}
	\hline
	\textbf{Látka} & \multicolumn{2}{|c|}{\textbf{pK$_S$}} &
	\multicolumn{2}{|c||}{\textbf{K$_S$}} &\textbf{Látka} & \multicolumn{2}{|c|}{\textbf{pK$_S$}} &
	\multicolumn{2}{|c|}{\textbf{K$_S$}} \\\hline
	AgBr & 12 & 31 & 4 & 90.10$^{-13}$ & CdS & 26 & 10 & 7 & 94.10$^{-27}$ \\\hline
	AgCl & 9 & 75 & 1 & 78.10$^{-10}$ & \ce{CoCO3} & 12 & 84 & 1 & 45.10$^{-13}$ \\\hline
	AgI & 16 & 08 & 8 & 32.10$^{-17}$ & CuBr & 8 & 28 & 5 & 25.10$^{-9}$ \\\hline
	AgSCN & 11 & 97 & 1 & 07.10$^{-12}$ & CuS & 35 & 20 & 6 & 31.10$^{-36}$ \\\hline
	\ce{Ag2CrO4} & 11 & 61 & 2 & 45.10$^{-12}$ & \ce{Ni(OH)2} & 15 & 50 & 3 & 16.10$^{-16}$  \\\hline
	\ce{Ag2S} & 49 & 20 & 6 & 31.10$^{-50}$ & PbS & 26 & 60 & 2 & 51.10$^{-27}$ \\\hline
	\ce{Ag2SO4} & 4 & 77 & 1 & 70.10$^{-5}$ & \ce{Sn(OH)4} & 56 & 00 & 1 & 00.10$^{-56}$ \\\hline
	\ce{BaCO3} & 8 & 29 & 5 & 13.10$^{-9}$ & \ce{Tl2S} & 20 & 30 & 5 & 01.10$^{-21}$ \\\hline
	\ce{BaCrO4} & 9 & 93 & 1 & 17.10$^{-10}$ & \ce{Zn3(AsO4)2} & 27 & 89 & 1 & 29.10$^{-28}$  \\\hline
\end{tabular}
\end{center}

\subsubsection{Koncentrace iontů}

\zadani{Vypočítejte koncentraci iontů v nasyceném roztoku AgCl}

Chlorid stříbrný v roztoku disociuje podle rovnice:

\ce{AgCl <=> Ag+ + Cl-}

Součin rozpustnosti vypočítáme:

$K_S = 10^{-9,75} = 1,78\times 10^{-10}\\
K_S = [\ce{Ag+}][\ce{Cl-}] = x^2\\
x = \sqrt{K_S} = \sqrt{1,78\times 10^{-10}} = 1,33\times 10^{-5}$ M

\odpoved{Koncentrace stříbrného a chloridového iontu jsou 1,33$\times$10$^{-5}$ mol.dm$^{-3}$}

\zadani{Vypočítejte koncentraci iontů v nasyceném roztoku \ce{Sn(OH)4}}

Hydroxid cíničitý, pKs = 56,00 , v roztoku disociuje podle rovnice:

\ce{Sn(OH)4 <=> Sn^{4+} + 4 OH-}

Součin rozpustnosti vypočítáme:

$K_S = 10^{-56,00} = 1 \times 10^{-56}\\
K_S = [\ce{Sn^{4+}}][\ce{OH-}]^4 = x \times (4x)^4 = 256x^5\\
x = \sqrt[5]{\frac{K_S}{256}} = \sqrt[5]{\frac{1 \times 10^{-56}}{256}} = 2,08\times 10^{-12}$ M = c(\ce{Sn^{4+}})

c(\ce{OH^{-}}) = 4 $\times\ 2,08\times10^{-12}$ = 8,33$\times10^{-12}$ M

\odpoved{Koncentrace cíničitých iontů je 2,08$\times10^{-12}$~mol.dm$^{-3}$ a hydroxidových iontů \\
je 8,33$\times$10$^{-12}$~mol.dm$^{-3}$.}

\clearpage